\section{Introduction}
\label{sec:carp:intro}

Organization of data generated by high-performance applications continues to be a challenge due to storage bandwidth constraints. Many high-performance applications consist of distinct computation phases, often called \emph{epochs} or \emph{timesteps}, that are periodically interrupted by checkpointing I/O phases. Examples of such applications include particle-in-cell frameworks~\cite{Byna2012}, mesh-based codes~\cite{Grete2023:Parthenon, Zhang2021:AMReX, Barke2024:Phoebus}, and machine learning model training. Periodic I/O phases enable scientists to perform queries with a temporal component, such as tracing trajectories of high-energy particles, or tracking shock waves in a fluid simulation. For efficient execution of such queries, persisted data is typically indexed in a post-processing stage. This post-processing requires additional time and resources, increasing the time to discovery. It also increases with the scale of the simulation, especially as computational bandwidth growth outpaces that of storage bandwidth~\cite{Bauer2016:SituMethodsInfrastructures}.

To reduce the time to discovery, in-situ approaches that index data as it is written by the parallel application have been proposed~\cite{Child2012:InSitu} as an alternative to post-processing. DeltaFS~\cite{Zheng2018:DeltaFS} intercepts application writes and organizes them into hash-based partitions. This enables point queries that retrieve individual items by key, but it breaks key locality. Range queries require key locality to efficiently find all items with keys in a given range. Range queries are an important tool for scientists interacting with continuous attributes such as temperature and velocity. They can be used to filter data into ranges of relevant values, such as isolating blast wave fronts or tracking energy bands. However, existing solutions for range query indexing force undesirable tradeoffs between write performance, query performance, and resource utilization (\S\ref{sec:carp:bg}).

To enable efficient range queries without post-processing, it is necessary to develop reordering-based in-situ indexing mechanisms. Reordering creates key locality in the data layout, which enables range queries to be served efficiently via a small number of contiguous writes. However, reordering of on-disk data stresses write performance by requiring I/Os for index maintenance. This leads to slower checkpointing and idling of compute nodes. Enabling efficient querying without adversely affecting write performance requires reordering data in the network data plane, before writing to storage. Doing so without requiring additional compute resources requires a lightweight embedded reorganization mechanism with a low memory footprint. These challenges are further compounded by the nature of scientific key distributions --- our analysis of real-world codes (\S\ref{sec:carp:distrib}) shows that scientific key distributions are highly skewed and vary over time. Reorganization mechanisms processing such data need to be adaptive to ensure that load hotspots in the reordering pipeline do not slow down application writes.

We present CARP (\emph{Continuous, Adaptive, Range Partitioner}), an adaptive in-situ indexing system that enables range queries without requiring data post-processing. CARP requires no changes to scientific applications and operates by intercepting writes in flight to storage. User input on key distribution characteristics is not required.  Partitions are instead discovered at runtime and adapted as the key distribution changes. CARP is write-optimized and does not divert any storage bandwidth from the application for reordering. Our evaluation shows that CARP incurs no overhead on application runtime and can achieve sub-second range query latency that is up to two orders of magnitude faster than state-of-the-art auxiliary indexing approaches. By partitioning data at runtime, CARP is up to 5$\times$ faster than post-processing approaches.

The contributions of this work are summarized as follows:

\begin{itemize}
  \item After a brief discussion on background (\S\ref{sec:carp:bg}), we characterize workload distributions and their drifts in modern scientific applications (\S\ref{sec:carp:distrib}) and demonstrate that static partitioning is inadequate for such workloads (\S\ref{sec:carp:eval_write}).
  \item We present CARP (\S\ref{sec:carp:des}), an in-situ data processing framework guaranteeing data locality and load balancing at scale. We identify requirements for a CARP storage backend and present KoiDB (\S\ref{sec:carp:des_store}), a reference backend implementation.
  \item We show that partial locality in data layout can be created in a single pass without requiring any reorganization and provide nearly-optimal range query performance (\S\ref{sec:carp:eval_read}), followed by a discussion on adapting CARP for different requirements (\S\ref{sec:carp:disc}).

\end{itemize}

We have released the source code~\cite{Jain2024:CARP} for both the CARP framework and the KoiDB backend.
